% ============================================================
%   Chapter 1  \textemdash\  Introduction
% ============================================================

\chapter{Introduction}\label{ch:introduction}

\section{Context and motivation}\label{sec:intro-context}

Road transport is one of the dominant causes of avoidable injury and death worldwide. According
to the World Health Organization, approximately 1.19 million people die every year as a direct
consequence of road traffic crashes, with human error being the leading contributing
factor~\cite{who2023road}. This figure has motivated a sustained effort, both in industry and in
academia, to automate vehicle control up to the point where the driver can be completely
displaced from the control loop. The SAE J3016 standard formalises this progression in six
levels, from no automation (Level 0) to full self-driving with no operational design domain
restrictions (Level 5)~\cite{sae2021j3016}. While most commercially deployed systems remain at
Level 2, research platforms increasingly target Levels 4 and 5, where the vehicle must cope with
the full variety of traffic situations, weather conditions and road topologies encountered in
real-world operation.

Among the techniques that have received growing attention for high-autonomy driving,
Reinforcement Learning (RL)~\cite{sutton2018rl} occupies a distinctive position. Classical
model-predictive and rule-based controllers require a hand-crafted description of every
situation and a correspondingly explicit control law; in contrast, an RL agent learns a policy
end-to-end by interacting with an environment and maximising a scalar reward signal. Recent
surveys of deep RL for autonomous driving report promising results on lane keeping, overtaking,
intersection handling and dense urban scenarios~\cite{kiran2022drlsurvey,aradi2022survey}.

The main obstacle to deploying these systems, however, is not their peak performance but their
behaviour \emph{during learning}. Vanilla RL relies on stochastic exploration, which is
incompatible with the safety requirements of a vehicle operating on public roads: even a single
exploratory collision is unacceptable. The research field generally referred to as
\emph{Safe Reinforcement Learning} (Safe RL) addresses this tension through multiple
strategies~\cite{garcia2015comprehensive}: encoding constraints in the reward function,
constraining the policy optimisation problem itself, or intercepting unsafe actions with a
runtime verifier known as a \emph{safety shield}~\cite{alshiekh2018shielding}. This project
belongs to the latter family: a PPO agent~\cite{schulman2017ppo} is trained with two alternative
shields that monitor its proposed actions, reason about their near-future consequences, and
substitute them when the resulting trajectory is judged unsafe. The shield layer runs entirely
during training, so safety violations are not merely discouraged by the reward function but
actively prevented at the actuator.

\section{Problem statement}\label{sec:intro-problem}

The central problem addressed in this thesis can be stated as follows. A vehicle, represented as
an RL agent, drives inside the CARLA simulator~\cite{dosovitskiy2017carla}. The agent observes a
partial representation of the environment through virtual sensors and a limited number of
kinematic variables, and must produce low-level control commands (steering, throttle and brake).
The goal is to learn a policy that

\begin{enumerate}
  \item \emph{maximises} a task reward that encodes progress, speed compliance and lane keeping,
  \item \emph{minimises} unsafe events such as collisions, off-road excursions and stalls, and
  \item \emph{tolerates} the shift in traffic density that occurs as the scenario is progressively
        complicated during training, and
  \item \emph{generalises} beyond the shield, so that the learned policy remains safe even when the
        shield is removed at deployment.
\end{enumerate}

The research question that guides the work is then: \emph{can a lightweight, interchangeable
safety-shield layer, combined with a dense reward shaper and a performance-driven curriculum,
train a PPO agent that is competitive with unconstrained PPO in terms of reward while keeping
crash and off-road rates close to zero during learning?} A second question, which emerges from the
first, is whether the very mechanism that enforces safety during training, the shield, induces
a \emph{dependence} that the agent cannot shed, and, if so, how that dependence can be measured and
reduced; this question is taken up in Chapter~\ref{ch:dependence}.

\section{Objectives}\label{sec:intro-objectives}

\subsection*{General objective}

To design, implement and empirically analyse a modular Safe RL framework for autonomous driving
on the CARLA simulator, combining a PPO agent with interchangeable safety shields, a dense
reward shaper and a curriculum manager.

\subsection*{Specific objectives}

\begin{itemize}
  \item[\textbf{SO1.}] Build a Gymnasium-compatible wrapper of the CARLA Python API that exposes a
        739-dimensional observation space (three semantic LIDAR channels plus lane-marking, lane,
        vehicle and route features) and a continuous two-dimensional action space, running in
        synchronous mode at 20\,Hz.
  \item[\textbf{SO2.}] Implement a PPO agent with generalised advantage estimation, clipped
        surrogate updates, value clipping, a target-KL early stop, cosine learning-rate annealing
        and online normalisation of the vector observation block.
  \item[\textbf{SO3.}] Develop a basic safety shield that combines LIDAR thresholds with
        information from CARLA's Waypoint API and applies a proportional correction to steering
        and brake whenever the agent's action is deemed unsafe.
  \item[\textbf{SO4.}] Develop an adaptive-horizon safety shield that classifies the current
        situation into three risk levels, predicts the vehicle's trajectory with a kinematic
        bicycle model, and selects the least invasive corrective action from a priority-ordered
        candidate set.
  \item[\textbf{SO5.}] Design a dense reward shaper that combines the sparse CARLA rewards with
        components for velocity tracking, lane centring, heading alignment, smoothness, progress,
        idle avoidance and lane-change tolerance.
  \item[\textbf{SO6.}] Build a performance-driven curriculum that progressively introduces
        non-playable vehicles (NPCs) into the scenario, with automatic rollback to earlier
        stages on performance collapse.
  \item[\textbf{SO7.}] Provide a live metrics pipeline based on SQLite (WAL mode) and a Streamlit
        dashboard so that each training run can be inspected while it is still in progress.
  \item[\textbf{SO8.}] Measure and reduce the agent's \emph{dependence} on the shield: instrument a
        shield-off evaluation probe, and study two strategies: a Large-Language-Model
        meta-controller that loosens the shield over training, and a behavioural-cloning scheme
        that turns the shield into a teacher, to train a policy that remains safe when the
        shield is removed at deployment.
\end{itemize}

\section{Scope and limitations}\label{sec:intro-scope}

The scope of this project is deliberately bounded to keep the contribution verifiable and
reproducible within the duration of a Bachelor's thesis.

\begin{itemize}
  \item The simulator version is \emph{CARLA 0.9.16}. Compatibility with earlier or future
        versions is not guaranteed.
  \item The main benchmark map is \emph{Town04}, a highway-oriented layout that is well suited
        for lane-keeping and overtaking experiments. Urban maps (Town01--Town03) and the maximum
        difficulty map (Town05) are supported by the framework but not systematically evaluated.
  \item Two shielding strategies are developed and compared: \emph{basic} (threshold + Waypoint
        API correction) and \emph{adaptive-horizon} (trajectory prediction with a kinematic
        bicycle model). Alternatives based on reachability analysis, model-predictive safety
        filters or formal verification are left for future work.
  \item The vehicle model used by the adaptive shield is the kinematic bicycle
        model~\cite{rajamani2012vehicle}. A dynamic model with tyre slip is not considered; at
        the speeds and lateral accelerations encountered in the experiments the kinematic
        approximation is a reasonable trade-off between fidelity and computational cost.
  \item Perception is based on semantic LIDAR plus CARLA's Waypoint API. Camera-based perception,
        sensor fusion and learned occupancy grids are not part of the observation space.
  \item The framework trains a \emph{single ego vehicle}. Multi-agent safe RL is an explicit
        direction for future extensions.
\end{itemize}

\section{Methodology overview}\label{sec:intro-methodology}

The project follows an iterative engineering methodology structured around four phases: design,
implementation, training and analysis, which are revisited in successive loops as new empirical
evidence becomes available.

\begin{enumerate}
  \item \textbf{Design.} The software architecture is expressed as a layered wrapper chain
        on top of Gymnasium. Each layer (environment, shield, reward shaper) has a single
        responsibility and a well-defined interface, so that components can be swapped or
        reconfigured without touching the others.
  \item \textbf{Implementation.} The Python code base is organised into self-contained modules
        (\texttt{src/CARLA/Env}, \texttt{src/PPO}, \texttt{src/Adaptative\_Shield},
        \texttt{src/reward\_shaper.py}, \texttt{src/curriculumManager.py}, \texttt{src/Metrics}).
        Automated style checks with Ruff and targeted unit tests on the reward shaper prevent
        regressions.
  \item \textbf{Training.} Each experimental configuration is launched via \texttt{main\_train.py}
        with explicit command-line flags. The agent writes metrics to a per-run SQLite database
        and checkpoints every $200$ episodes.
  \item \textbf{Analysis.} The \texttt{dataVisualizer.py} Streamlit dashboard is used to inspect
        reward curves, safety rates, curriculum transitions and shield activations during and
        after training. Observed pathologies are translated into new design iterations.
\end{enumerate}

\section{Document structure}\label{sec:intro-structure}

The remainder of this document is organised as follows.

\begin{description}
  \item[Chapter~\ref{ch:sota}] reviews the theoretical background. It introduces Markov Decision
        Processes, derives the PPO objective from the policy gradient, describes the Generalised
        Advantage Estimator, surveys the main families of Safe RL, formalises the kinematic
        bicycle model and summarises the features of the CARLA simulator that are relevant to
        this project.
  \item[Chapter~\ref{ch:requirements}] captures the functional and non-functional requirements
        that the framework must satisfy, together with the technical constraints imposed by the
        simulator, the hardware and the Bachelor-thesis timeline.
  \item[Chapter~\ref{ch:design}] describes the system design: the wrapper chain, the structure
        of the observation and action spaces, the Actor\textendash Critic network, the two
        safety shields, the reward shaper and the curriculum manager.
  \item[Chapter~\ref{ch:implementation}] details the implementation: technology stack, project
        layout, key algorithms in pseudocode, training and evaluation pipelines, and
        reproducibility considerations.
  \item[Chapter~\ref{ch:experimentation}] presents the experimental protocol and a structured
        template to report the empirical results obtained on Town04.
  \item[Chapter~\ref{ch:dependence}] studies the agent's dependence on the shield: it introduces
        the shield-off evaluation probe, reports a negative result for a Large-Language-Model
        meta-controller that loosens the shield, and a positive result for a behavioural-cloning
        scheme that turns the shield into a teacher of unaided safe driving.
  \item[Chapter~\ref{ch:conclusions}] summarises the contributions, evaluates the fulfilment of
        the objectives enumerated above, and proposes future work.
  \item[Appendices~\ref{app:hyperparameters}, \ref{app:observation} and \ref{app:reproducibility}]
        provide the complete list of default hyperparameters, a full breakdown of the observation
        space and a reproducibility guide, respectively.
\end{description}
